\documentclass[12pt]{article}
\usepackage[utf8]{inputenc}
\usepackage[T1]{fontenc}
\usepackage[margin=1in]{geometry}
\usepackage{setspace}
\usepackage{xcolor}
\usepackage[colorlinks=true, urlcolor=blue!60!black, linkcolor=blue!60!black]{hyperref}
\setlength{\parindent}{2em}
\setlength{\parskip}{1.25em}
\renewcommand{\baselinestretch}{1.0}
\begin{document}
\begin{center}
{\LARGE \textit{The Prize Cases}, 67 U.S. 635 (1863)\par}
\vspace{0.8em}
{\large Case Brief\par}
\end{center}
\vspace{1em}

\section{Memory Jogger}

In a 5--4 wartime decision, the Supreme Court upheld President Lincoln's unilateral April 1861 naval blockade of Confederate ports, ruling that while the President cannot initiate or declare war, he is constitutionally bound to repel sudden invasions or armed rebellions by force without waiting for a congressional declaration of war (\href{https://www.courtlistener.com/opinion/87541/the-brig-amy-warwick/}{CourtListener}; \href{https://supreme.justia.com/cases/federal/us/67/635/}{Justia}).

\section{Detailed Case Facts}

Following the April 12, 1861 Confederate bombardment of Fort Sumter---while Congress was out of session---President Abraham Lincoln issued a series of emergency proclamations invoking the Militia Act of 1795 (1 Stat. 424) and the Insurrection Act of 1807 (2 Stat. 443) (\href{https://www.law.cornell.edu/constitution-conan/article-1/section-8/clause-1/civil-war-war-powers-and-the-prize-cases}{Constitution Annotated}; \href{https://supreme.justia.com/cases/federal/us/67/635/}{Justia}). On April 19, 1861, Lincoln proclaimed a naval blockade of ports in South Carolina, Georgia, Alabama, Florida, Mississippi, Louisiana, and Texas, and on April 27, 1861, he extended the blockade to Virginia and North Carolina (\href{https://supreme.justia.com/cases/federal/us/67/635/}{Justia}). Commodore Pendergrast announced the Virginia blockade as established on April 30, 1861, with fifteen days allowed for neutral vessels to depart (\href{https://supreme.justia.com/cases/federal/us/67/635/}{Justia}). When Congress convened in special session in July 1861, it enacted the Act of July 13, 1861 (12 Stat. 257), authorizing the President to declare states in insurrection and interdict trade, and the Act of August 6, 1861 (12 Stat. 326), retroactively approving, legalizing, and making valid Lincoln's prior military acts and proclamations (\href{https://www.law.cornell.edu/constitution-conan/article-1/section-8/clause-1/civil-war-war-powers-and-the-prize-cases}{Constitution Annotated}; \href{https://www.oyez.org/cases/1850-1900/67us635}{Oyez}).

Between May and July 1861---all prior to the July 13, 1861 Act---Union naval vessels captured four merchant ships and their cargoes as prizes of war (\href{https://supreme.justia.com/cases/federal/us/67/635/}{Justia}):
\begin{itemize}
\item \textit{The Brig Amy Warwick}: A Richmond-owned American merchant vessel carrying 5,100 bags of coffee from Rio de Janeiro to Hampton Roads (for delivery at New York, Philadelphia, Baltimore, or Richmond), captured off Cape Henry on July 10, 1861, by the USS \textit{Quaker City} while its master was ignorant of the war, and libeled as enemy property in Boston (\href{https://supreme.justia.com/cases/federal/us/67/635/}{Justia}).
\item \textit{The Schooner Crenshaw}: A Richmond-owned vessel carrying tobacco from Richmond to Liverpool under a Confederate port clearance, captured in Hampton Roads at the mouth of the James River on May 17, 1861, by the USS \textit{Star} and taken to New York (\href{https://supreme.justia.com/cases/federal/us/67/635/}{Justia}).
\item \textit{The Barque Hiawatha}: A British-owned vessel carrying tobacco and cotton from Richmond to Liverpool; it finished loading on May 15, 1861, was delayed by lack of a steam-tug until May 17--18, and was captured in Hampton Roads on May 20, 1861, by the USS \textit{Minnesota} and taken to New York (\href{https://supreme.justia.com/cases/federal/us/67/635/}{Justia}).
\item \textit{The Schooner Brilliante}: A Mexican-owned vessel carrying flour from New Orleans to Sisal and Campeche; after being warned off the Mississippi River blockade by the USS \textit{Brooklyn}, it entered New Orleans via Lake Pontchartrain and was captured on June 23, 1861, in Biloxi Bay by boats from the USS \textit{Massachusetts} and taken to Key West (\href{https://supreme.justia.com/cases/federal/us/67/635/}{Justia}).
\end{itemize}

\section{Procedural History}

Federal district attorneys filed prize libels on behalf of the United States and the capturing naval officers and crews in the respective federal district courts sitting in admiralty: the District of Massachusetts (\textit{Amy Warwick}), the Southern District of New York (\textit{Crenshaw} and \textit{Hiawatha}), and the Southern District of Florida (\textit{Brilliante}) (\href{https://supreme.justia.com/cases/federal/us/67/635/}{Justia}; \href{https://caselaw.findlaw.com/court/us-supreme-court/67/635.html}{FindLaw}). In each case, the district court entered a decree of condemnation (restoring certain neutral or loyal-citizen cargo interests in \textit{Amy Warwick} and \textit{Crenshaw}) (\href{https://supreme.justia.com/cases/federal/us/67/635/}{Justia}). The ship and cargo claimants appealed directly to the Supreme Court, arguing that no legal war existed prior to July 13, 1861, and that the President lacked constitutional authority to institute a blockade without a congressional declaration of war (\href{https://supreme.justia.com/cases/federal/us/67/635/}{Justia}). The consolidated appeals were argued over twelve days from February 10 through February 25, 1863, and decided on March 10, 1863 (\href{https://supreme.justia.com/cases/federal/us/67/635/}{Justia}).

\section{Judicial Votes}

The Court decided the consolidated cases (reported at 67 U.S. (2 Black) 635) by a \textbf{5--4} vote (\href{https://teachingamericanhistory.org/document/the-prize-cases/}{Teaching American History}; \href{https://supreme.justia.com/cases/federal/us/67/635/}{Justia}):
\begin{itemize}
\item \textbf{Majority (5):} Justice Robert Cooper Grier (author), joined by Justices James Moore Wayne, Noah Haynes Swayne, Samuel Freeman Miller, and David Davis (\href{https://teachingamericanhistory.org/document/the-prize-cases/}{Teaching American History}).
\item \textbf{Concurrences (0):} None.
\item \textbf{Dissent (4):} Justice Samuel Nelson (author), joined by Chief Justice Roger B. Taney and Justices John Catron and Nathan Clifford (\href{https://supreme.justia.com/cases/federal/us/67/635/}{Justia}).
\end{itemize}

\section{Holding}

\textbf{Judgment:} The decrees of condemnation were affirmed with costs in \textit{The Amy Warwick}, \textit{The Hiawatha}, and \textit{The Brilliante}. In \textit{The Crenshaw}, the condemnation of the vessel and enemy-owned cargo was affirmed, the costs decree against neutral claimants Lear \& Sons was affirmed, and thirty tierces of tobacco strips purchased before the war by New York citizens Irvin \& Co. were ordered restored (\href{https://supreme.justia.com/cases/federal/us/67/635/}{Justia}).

\textbf{Rule Announced:} First, under Article II and the Acts of 1795 and 1807, although the President has no constitutional power to initiate or declare war, when a foreign invasion or armed domestic rebellion of civil-war magnitude occurs, the President is both authorized and bound to resist force by force---including by instituting a naval blockade \textit{jure belli}---without waiting for a formal declaration of war by Congress (\href{https://supreme.justia.com/cases/federal/us/67/635/}{Justia}; \href{https://www.law.cornell.edu/constitution-conan/article-1/section-8/clause-1/civil-war-war-powers-and-the-prize-cases}{Constitution Annotated}). Second, the President's determination that an insurrection has reached the proportions of a civil war warranting belligerent measures is a political judgment to which courts must defer (\href{https://supreme.justia.com/cases/federal/us/67/635/}{Justia}). Third, property belonging to persons domiciled in hostile territory held by the Confederacy is subject to maritime capture as ``enemies' property'' regardless of the owner's personal loyalty (\href{https://supreme.justia.com/cases/federal/us/67/635/}{Justia}).

\textbf{Holding vs. Dicta:} The core holding establishes that President Lincoln lawfully instituted the April 1861 blockade and captured ships as prizes prior to July 13, 1861, and that residents of Confederate-held territory could be treated as enemies for maritime prize purposes. Justice Grier's alternative observation that Congress's August 6, 1861 ratification statute ``operated to perfectly cure the defect'' if any legislative sanction had been needed (``[w]ithout admitting that such an act was necessary'') is an alternative ground or dicta (\href{https://supreme.justia.com/cases/federal/us/67/635/}{Justia}; \href{https://www.law.cornell.edu/constitution-conan/article-2/section-2/clause-1/prize-cases-and-commander-in-chief-clause}{Constitution Annotated}). Likewise, Grier's statement that Congress ``cannot declare war against a State, or any number of States'' was a broad constitutional gloss explaining why civil wars begin de facto rather than by formal declaration (\href{https://supreme.justia.com/cases/federal/us/67/635/}{Justia}).

\section{Analysis of Opinions}

\subsection{Majority Opinion (Justice Grier)}

Justice Grier framed two overarching questions: (1) whether the President had authority under international and constitutional law to institute a blockade of rebellious ports, and (2) whether property of residents of seceding states could be captured at sea as ``enemies' property'' (\href{https://supreme.justia.com/cases/federal/us/67/635/}{Justia}). On the first question, Grier drew on the law of nations (citing Vattel and Lord Stowell) to explain that war is a matter of fact, not mere legal formality: a civil war is ``never solemnly declared,'' but exists whenever armed rebels occupy territory, establish a de facto government, organize armies, and close the regular courts of justice (\href{https://supreme.justia.com/cases/federal/us/67/635/}{Justia}). While Congress alone holds the Article I power to declare a national or foreign war and the President ``has no power to initiate or declare a war,'' the Constitution vests the President with the Executive Power, the Take Care Clause, and Commander-in-Chief authority, reinforced by the Acts of 1795 and 1807 (\href{https://supreme.justia.com/cases/federal/us/67/635/}{Justia}). When war is thrust upon the nation by invasion or rebellion, the President is ``bound to meet it in the shape it presented itself, without waiting for Congress to baptize it with a name'' (\href{https://supreme.justia.com/cases/federal/us/67/635/}{Justia}).

Grier further held that the judiciary must defer to the President's political judgment on whether an insurrection has escalated to belligerent war, noting that Queen Victoria's May 13, 1861 neutrality proclamation estopped foreign neutrals from denying that a war existed (\href{https://supreme.justia.com/cases/federal/us/67/635/}{Justia}). Even if prior congressional approval had been required, Grier added, Congress's August 6, 1861 Act retroactively ratified Lincoln's proclamations, and the Ex Post Facto Clause did not bar retroactive validation in admiralty prize courts (\href{https://supreme.justia.com/cases/federal/us/67/635/}{Justia}). On the second question, Grier rejected the claimants' argument that Southern residents were merely domestic traitors immune from prize capture: in a territorial civil war bounded by ``lines of bayonets,'' a sovereign may exercise both sovereign and belligerent rights, and maritime commerce originating in hostile territory is subject to capture as enemies' property without regard to individual allegiance (\href{https://supreme.justia.com/cases/federal/us/67/635/}{Justia}).

\subsection{Concurring Opinions}

No concurring opinions were filed (\href{https://supreme.justia.com/cases/federal/us/67/635/}{Justia}).

\subsection{Dissenting Opinion (Justice Nelson)}

Justice Nelson (joined by Chief Justice Taney and Justices Catron and Clifford) agreed that Congress had lawfully recognized a territorial civil war in the Act of July 13, 1861, but insisted that all captures made \textit{before} July 13, 1861, were illegal and void (\href{https://supreme.justia.com/cases/federal/us/67/635/}{Justia}). Nelson drew a sharp line between war ``in a material sense'' and war ``in a legal sense'' under the Constitution and the law of nations (\href{https://supreme.justia.com/cases/federal/us/67/635/}{Justia}). Because a legal state of war transforms peaceful citizens into public enemies, suspends contracts, voids commerce, subjects neutral shipping to capture, and authorizes confiscation without criminal trial, the Constitution vests the war-making power exclusively in Congress under Article I, Section 8, Clause 11 (\href{https://supreme.justia.com/cases/federal/us/67/635/}{Justia}).

Nelson rejected the claim that necessity compelled unilateral presidential war powers: under the Acts of 1795 and 1807, Lincoln already possessed full municipal authority to deploy the entire army and navy in a ``personal war'' against armed insurgents while Congress was summoned within thirty days (\href{https://supreme.justia.com/cases/federal/us/67/635/}{Justia}). Analogizing to Great Britain's handling of the American Revolution---which remained a personal conflict until Parliament's Prohibitory Act of 1776 made it a territorial war---Nelson argued that the President could not unilaterally convert loyal citizens into belligerent enemies (\href{https://supreme.justia.com/cases/federal/us/67/635/}{Justia}). He also rejected retroactive congressional ratification under the August 1861 Act as an impermissible ex post facto confiscation of property lawfully engaged in trade at the time of seizure (\href{https://supreme.justia.com/cases/federal/us/67/635/}{Justia}). On the specific facts of \textit{The Hiawatha}, Nelson additionally concluded that the vessel had no intent to break the blockade and was entitled under the text of Lincoln's April 19 proclamation to an endorsed warning before seizure (\href{https://supreme.justia.com/cases/federal/us/67/635/}{Justia}).

\section{Examples: Future Applications}

\begin{enumerate}
\item \textbf{Same-Side Application (Repelling a Sudden Foreign Attack Without Prior Congressional Authorization):} While Congress is in recess, a foreign naval fleet launches a surprise missile strike on a U.S. naval base in the Pacific; the President immediately orders U.S. carrier groups to destroy the attacking fleet and establish a maritime exclusion zone around the hostile staging port. \textit{Analysis:} Under \textit{The Prize Cases}, the President's defensive response is constitutionally valid. The President did not initiate the war but is ``authorized [and] bound to resist force by force'' to meet a sudden attack without waiting for a formal declaration of war (\href{https://supreme.justia.com/cases/federal/us/67/635/}{Justia}).

\item \textbf{Same-Side Application (Suppressing a Large-Scale Armed Insurrection):} An organized armed separatist movement seizes federal forts, expels federal courts across two states, and fields uniformed regiments; the President invokes the Insurrection Act and imposes a naval and air quarantine on the insurgent-held territory before Congress can assemble. \textit{Analysis:} This falls squarely within \textit{The Prize Cases}: when armed rebellion interrupts the regular course of justice so that courts cannot remain open, the President may meet the de facto civil war with belligerent defensive measures, and courts will defer to the President's assessment of the force required (\href{https://supreme.justia.com/cases/federal/us/67/635/}{Justia}).

\item \textbf{Opposite-Side Application (Unilateral Preventive War Against a Non-Attacking Foreign State):} Intelligence suggests a hostile foreign nation may develop advanced weapons within five years; although the nation has not attacked the United States or U.S. forces, the President orders a full-scale air and ground invasion and naval blockade without seeking congressional authorization. \textit{Analysis:} Under \textit{The Prize Cases}, this exceeds presidential power. Justice Grier explicitly affirmed that the President ``has no power to initiate or declare a war either against a foreign nation or a domestic State,'' limiting unilateral executive force to defensive responses to actual attack or armed rebellion (\href{https://supreme.justia.com/cases/federal/us/67/635/}{Justia}).

\item \textbf{Opposite-Side Application (Domestic Seizure of Private Industry or Military Trials Where Civil Courts Are Open):} During an overseas conflict, the President unilaterally seizes domestic steel mills to avert a labor strike (without statutory authorization) or tries civilian citizens in military tribunals in a loyal state where federal courts are open and functioning. \textit{Analysis:} Unlike maritime prize captures in an active theater of rebellion, unilateral domestic seizures and military trials of civilians outside the war zone exceed Commander-in-Chief authority, as the Court later held in \href{https://sites.gsu.edu/us-constipedia/ex-parte-milligan-1866/}{\textit{Ex parte Milligan}, 71 U.S. 2 (1866)}, and \href{https://supreme.justia.com/cases/federal/us/343/579/}{\textit{Youngstown Sheet \& Tube Co. v. Sawyer}, 343 U.S. 579 (1952)}.

\item \textbf{Fence-Sitter (Unilateral Presidential Military Strikes in Response to Cyberattacks or Attacks on Foreign Allies):} A hostile state launches a severe, state-sponsored cyberattack that disables a U.S. regional power grid---or fires missiles at a treaty ally's commercial vessels in international waters---and the President orders retaliatory kinetic airstrikes and a limited naval interdiction without congressional approval. \textit{Analysis:} This sits on the constitutional boundary of \textit{The Prize Cases}. Proponents argue that a crippling cyberattack or attack on vital national interests is the modern equivalent of a sudden invasion requiring immediate defensive force; critics counter that absent a physical attack on U.S. territory or forces, escalating to kinetic hostilities amounts to initiating war, which Article I reserves to Congress (\href{https://www.law.cornell.edu/constitution-conan/article-1/section-8/clause-1/civil-war-war-powers-and-the-prize-cases}{Constitution Annotated}).
\end{enumerate}

\section{Critique}

\textbf{Originalist and Textualist Perspectives:} Originalist war-powers scholars are deeply divided over \textit{The Prize Cases}. On one hand, originalists such as Michael D. Ramsey and Saikrishna Prakash (along with Founding-era analyses of the Declare War Clause) emphasize that the Framers deliberately transferred the war-initiation power from the executive to Congress, leaving the President only a narrow defensive authority to ``repel sudden attacks''---making Grier's explicit statement that the President ``has no power to initiate or declare a war'' faithful to original meaning, even as they caution against reading the decision to authorize offensive unilateral wars (\href{https://southerncalifornialawreview.com/2024/02/20/delegating-war-powers/}{Southern California Law Review}; \href{https://lawreview.uchicago.edu/sites/default/files/78-1-ExecutivePower\%20and\%20the\%20Discipline\%20of\%20History-Mortenson.pdf}{U. Chi. Law Review}). Conversely, executive-power theorists such as John C. Yoo argue that Article II's Vesting and Commander-in-Chief Clauses confer broad independent authority to initiate and direct military engagements unless Congress uses its funding or legislative powers to check the President (\href{https://www.law.berkeley.edu/archive/files/yoowarpowers.pdf}{Yoo, UC Berkeley Law}). Other scholars, examining the Calling Forth Clause (Article I, Section 8, Clause 15), argue that the Court's heavy reliance on inherent Article II power obscured the Constitution's textual design placing domestic insurrection responses under congressional statutory control (\href{https://constitutioncenter.org/the-constitution/articles/article-ii/clauses/345}{National Constitution Center}; \href{https://journals.law.harvard.edu/nsj/wp-content/uploads/sites/82/2014/01/Bahar-Presidential-Intervention-Principle.pdf}{Bahar, Harvard National Security Journal}).

\textbf{Doctrinal and Separation-of-Powers Perspectives:} Constitutional war-powers scholars such as Louis Fisher and David J. Barron \& Martin S. Lederman argue that modern administrations have routinely overread \textit{The Prize Cases} as a blank check for unilateral presidential warmaking (\href{https://scholarship.law.slu.edu/cgi/viewcontent.cgi?article=1581&context=lj}{Pearlstein, SLU Law Journal}; \href{https://harvardlawreview.org/wp-content/uploads/2007/12/barron_lederman.pdf}{Barron \& Lederman, Harvard Law Review}). They emphasize two crucial limits built into the decision itself: Lincoln acted defensively against an actual domestic rebellion while Congress was out of session, and Congress swiftly ratified his measures by statute (\href{https://www.law.cornell.edu/constitution-conan/article-2/section-2/clause-1/prize-cases-and-commander-in-chief-clause}{Constitution Annotated}). Deborah Pearlstein likewise notes that despite strong executive pressure for judicial abstention, the Court refused to abdicate its role entirely, scrutinizing individual captures under international prize law and restoring unlawfully condemned property in \textit{The Crenshaw} (\href{https://scholarship.law.slu.edu/cgi/viewcontent.cgi?article=1581&context=lj}{Pearlstein, SLU Law Journal}).

\textbf{Underlying Competing Values:} The decision reflects the perennial constitutional tension between \textit{executive energy and national survival} (enabling decisive, immediate defense when the Republic faces existential crisis before Congress can assemble) and \textit{republican deliberation and the rule of law} (requiring collective legislative consent before the nation crosses the threshold from peace to war and strips citizens and neutrals of peacetime legal protections).

\section{Key Quotations}

\begin{quote}
``He has no power to initiate or declare a war either against a foreign nation or a domestic State. \ldots If a war be made by invasion of a foreign nation, the President is not only authorized but bound to resist force by force. He does not initiate the war, but is bound to accept the challenge without waiting for any special legislative authority.''
\end{quote}
\noindent --- Justice Robert Cooper Grier, Opinion of the Court, 67 U.S. at 668 (\href{https://supreme.justia.com/cases/federal/us/67/635/}{Justia}; \href{https://www.courtlistener.com/opinion/87541/the-brig-amy-warwick/}{CourtListener}).

\begin{quote}
``The President was bound to meet it in the shape it presented itself, without waiting for Congress to baptize it with a name; and no name given to it by him or them could change the fact.''
\end{quote}
\noindent --- Justice Robert Cooper Grier, Opinion of the Court, 67 U.S. at 669 (\href{https://supreme.justia.com/cases/federal/us/67/635/}{Justia}).

\begin{quote}
``Whether the President, in fulfilling his duties as Commander-in-chief in suppressing an insurrection, has met with such armed hostile resistance and a civil war of such alarming proportions as will compel him to accord to them the character of belligerents is a question to be decided by him, and this Court must be governed by the decisions and acts of the political department of the Government to which this power was entrusted.''
\end{quote}
\noindent --- Justice Robert Cooper Grier, Opinion of the Court, 67 U.S. at 670 (\href{https://supreme.justia.com/cases/federal/us/67/635/}{Justia}).

\begin{quote}
``For we find there that to constitute a civil war in the sense in which we are speaking, before it can exist in contemplation of law, it must be recognized or declared by the sovereign power of the State, and which sovereign power by our Constitution is lodged in the Congress of the United States -- civil war, therefore, under our system of government, can exist only by an act of Congress, which requires the assent of two of the great departments of the Government, the Executive and Legislative.''
\end{quote}
\noindent --- Justice Samuel Nelson, Dissenting Opinion, 67 U.S. at 690 (\href{https://supreme.justia.com/cases/federal/us/67/635/}{Justia}).

\begin{quote}
``I am compelled to the conclusion that no civil war existed between this Government and the States in insurrection till recognized by the Act of Congress 13th of July, 1861; that the President does not possess the power under the Constitution to declare war or recognize its existence within the meaning of the law of nations, which carries with it belligerent rights, and thus change the country and all its citizens from a state of peace to a state of war; that this power belongs exclusively to the Congress of the United States''
\end{quote}
\noindent --- Justice Samuel Nelson, Dissenting Opinion, 67 U.S. at 698--99 (\href{https://supreme.justia.com/cases/federal/us/67/635/}{Justia}).

\end{document}
